\chapter{Introduzione}\label{chap:first}

\inlineminitoc

\section{Definizione del problema}
Il lavoro di tesi si inserisce nel contesto delle attività di ricerca e sviluppo svolte presso InfoCube, un'azienda di consulenza informatica specializzata, tra le altre cose, in Intelligenza Artificiale.
Nello specifico, il progetto ruota attorno a \textbf{LUDOVICA} (\textit{Language Understanding for Document Optimization and Validation In Civil Administration}), un motore generativo basato sull'IA progettato per la Pubblica Amministrazione Locale.
L'obiettivo principale di LUDOVICA è fornire ai funzionari un sistema capace di elaborare richieste in linguaggio naturale, cercare atti pertinenti, e generare template di atti amministrativi pronti all'uso e modificabili.

Il problema tecnico fondamentale affrontato in questo tirocinio consiste nella \textbf{trasformazione di documenti legali in un formato strutturato e interrogabile dalle macchine, in modo da poter costruire un solido database a grafo (\textit{Knowledge Graph})}.
Poiché gli atti normativi possiedono una gerarchia fissa e complessa (Titoli, Capi, Articoli, Commi), è cruciale che l'architettura dei dati rifletta questa struttura per permettere al sistema di risalire sempre al documento originale a partire da una singola porzione di testo.

A livello operativo, le attività di tirocinio si sono concentrate su tre task principali:
\begin{enumerate}
    \item Il \textbf{salvataggio dei documenti legali e delle loro gerarchie} sul database a grafo \textit{Neo4j}.
    \item L'\textbf{implementazione della \textit{Semantic Search}}, affiancata al \textbf{testing di molteplici modelli di embedding} per valutarne le prestazioni e l'accuratezza delle risposte sul testo legale in lingua italiana.
    \item L'\textbf{introduzione del database vettoriale\textit{Qdrant}}, impiegato unicamente a scopo di testing per poter effettuare un \textbf{confronto diretto delle performance} tra un approccio basato sulla sola ricerca vettoriale e quello topologico basato sui grafi.
\end{enumerate}

\section{Tecnologie utilizzate}
Per far fronte alle sfide di estrazione, elaborazione, archiviazione e interrogazione dei documenti normativi, è stato impiegato uno stack tecnologico che combina strumenti di parsing avanzato, modelli di LLM ed embedding e database specializzati:

\subsection*{Linguaggi di Programmazione}
\begin{itemize}
    \item \textbf{Python:} Scelto per la vastità e la potenza delle sue librerie, è stato il linguaggio cardine per la costruzione delle pipeline del progetto.
\end{itemize}

\subsection*{Strumenti per l'Estrazione del Testo}
\begin{itemize}
   \item \textbf{Open Data Loader:} Una libreria open-source scelta per la sua capacità di estrarre testo in modo \textit{deterministico} da file PDF, convertendolo direttamente in Markdown.
    \item \textbf{Docling:} Una libreria open-source di parsing e analisi del layout documentale, specializzato nell'estrazione precisa di testo, tabelle e metadati strutturati a partire da formati complessi come i PDF. Integrato come supporto di Open Data Loader.
    \item \textbf{RapidOCR:} Modulo impiegato per estrarre testo dai documenti derivanti da scansioni visive.
    Tra i vari motori OCR disponibili, è stato optato l'uso di questo specifico OCR che si è rivelato il compromesso più efficace per bilanciare i tempi di esecuzione e l'elevato consumo di RAM del container.
\end{itemize}

\subsection*{Database e Linguaggi di Interrogazione}
\begin{itemize}
    \item \textbf{Neo4j:} Il database a grafo principale utilizzato per mappare la rigida struttura gerarchica degli atti normativi e i nodi concettuali.
Consente di gestire e navigare le relazioni in maniera altamente efficiente rispetto ai database relazionali tradizionali.
    \item \textbf{Cypher:} Il linguaggio di interrogazione dichiarativo nativo di Neo4j, utilizzato per effettuare le query di inserimento (ingestion) e di ricerca topologica nel grafo.
    \item \textbf{Qdrant:} Un database vettoriale utilizzato sperimentalmente per condurre test e confronti sulle performance della ricerca semantica (\textit{Cosine Similarity}) rispetto all'approccio basato sui grafi.
\end{itemize}

\subsection*{LLM e Modelli di Embedding}
\begin{itemize}
    \item \textbf{Gemini 2.5-flash:} \textit{Large Language Model} utilizzato tramite API per il compito cruciale della strutturazione sintattica.
    Riceveva in input il testo in formato Markdown dai PDF e restituiva JSON formattati rigorosamente secondo uno schema predefinito, che rappresentavano la gerarchia dei documenti normativi (Titoli, Capi, Articoli, Commi) e i nodi concettuali.
    \item \textbf{Modelli di Embedding:} Sono stati testati svariati modelli per ottimizzare la vettorializzazione del testo legale in italiano. I test più approfonditi hanno coinvolto i seguenti modelli: \texttt{intfloat/multilingual-e5-large}, \texttt{google/gemini-embedding-001} e \texttt{BAAI/bge-m3}. La valutazione si è basata sulla capacità di recuperare commi pertinenti a partire da query in linguaggio naturale.
\end{itemize}


\subsection*{Deployment e Containerizzazione}
\begin{itemize}
    \item \textbf{Docker:} L'intero ecosistema è stato modularizzato attraverso lo sviluppo di container isolati e specifici (es. \texttt{extract-from-pdf}, \texttt{db-ingestion}, \ldots{}). Questa scelta ha garantito un'elevata modularità, l'isolamento delle dipendenze software e un rigoroso controllo sulle risorse computazionali, specialmente durante i task più onerosi legati all'\textbf{\textit{Optical Character Recognition (OCR)}} e alla \textbf{\textit{generazione degli Embedding}}.
\end{itemize}